\section{Introduction}

The hardware a student actually owns is rarely the hardware the tooling assumes.
Over a few years a person may accumulate two or three laptops, each adequate on
its own terms, none of them carrying a large GPU. The machines used throughout
this report are a case in point: a Linux laptop with an RTX~3050 exposing
\SI{4096}{MiB} of video memory, a Windows laptop with an RTX~4050 exposing
\SI{6141}{MiB}, and an M4 MacBook~Air with \SI{16384}{MiB} of unified memory.
Individually, none of them holds a 4.9~GB quantised model in GPU memory.
Together they hold more than enough.

\texttt{llama.cpp} already provides the mechanism for combining them. Its RPC
backend exposes a remote accelerator over TCP, and the main process can offload
model layers to it. The upstream documentation is explicit that this backend is
a proof of concept and that it should not be exposed to untrusted networks. It
is equally explicit in behaviour: if an RPC peer stops responding, the process
using it aborts. We reproduced this directly, and the error text is quoted in
Section~\ref{sec:background}.

That behaviour is the obstacle. A laptop is not a server. Its lid gets closed,
its Wi-Fi roams, its owner stops a terminal. Each of those events ends the
inference process outright, and the model must be loaded again from the
beginning. On the network used here, reloading a 3B model across two machines
took between \SI{33}{\second} and \SI{240}{\second} depending on cache state,
so an unhandled failure is not a brief interruption.

\subsection{What this report covers}

HyperPod is the control plane we built around \texttt{llama.cpp} to make that
pool usable. It does not implement attention kernels, quantisation, or device
offload; \texttt{llama.cpp} performs all of that. HyperPod is responsible for
five things:

\begin{enumerate}[leftmargin=1.4em,itemsep=2pt]
  \item finding worker laptops on the local network without manual addresses;
  \item recording what each laptop reports about its own hardware, and marking
        anything it cannot report as unavailable rather than guessing;
  \item deriving tensor weights from measured free video memory;
  \item reading back the placement \texttt{llama.cpp} actually chose, rather
        than assuming the requested split was honoured;
  \item detecting that the model process has died, excluding the endpoint
        responsible, rebuilding on what remains, and retrying the request that
        was interrupted.
\end{enumerate}

Section~\ref{sec:design} describes the design, including the one scheduling
decision that required a measurement to settle. Section~\ref{sec:results}
reports results from the three machines above: the placement
\texttt{llama.cpp} produced, five-run latency and throughput figures against
two baselines, and two physical failure experiments in which an accelerator was
stopped by hand during a live request.

\subsection{What this report does not claim}

Capability-aware placement across heterogeneous devices is not new, and neither
is failure recovery in distributed serving. Both exist in datacentre systems and
in prior work on consumer-hardware inference, discussed in
Section~\ref{sec:related}. The contribution here is narrower: a working,
measured implementation on three mismatched laptops across CUDA and Metal, with
every claim backed by a committed log, and with the costs reported rather than
omitted.

Three specific costs are reported in full. Splitting over Wi-Fi was roughly six
times slower per output token than the single local GPU. Recovery restarts the
model process and re-runs the request; it does not resume the KV cache or
continue from the token where the failure occurred. And the 8B configuration
demonstrates pooled capacity only: because that model does not fit in the
controller's GPU memory alone, the two-machine arrangement cannot survive the
loss of the remote worker. Each of these is stated again in
Section~\ref{sec:limitations} so that no reader has to infer them.
